\documentclass[11pt,a4paper]{article}
\usepackage[T1]{fontenc}
\usepackage[utf8]{inputenc}
\usepackage{amsmath,amsthm,mathtools}
\usepackage{newpxtext,newpxmath}
\usepackage[a4paper,margin=25mm,headheight=15pt]{geometry}
\usepackage{microtype}
\usepackage{xcolor,booktabs,array,longtable}
\usepackage{enumitem}
\usepackage{fancyhdr}
\usepackage{listings}
\usepackage[most]{tcolorbox}
\usepackage{hyperref}
\definecolor{ink}{HTML}{183547}
\definecolor{accent}{HTML}{126B73}
\definecolor{pale}{HTML}{EEF5F6}
\hypersetup{colorlinks=true,linkcolor=accent,citecolor=accent,urlcolor=accent,
  pdftitle={Exact maximal order types below omega squared},
  pdfauthor={ChatGPT, research draft prepared for Vladimir Reshetnikov},
  pdfsubject={A proposed solution to Altman's finite-alphabet sharpness question}}
\pagestyle{fancy}
\fancyhf{}
\fancyhead[L]{\small\textsc{Exact order types below $\omega^2$}}
\fancyhead[R]{\small Research draft}
\fancyfoot[C]{\thepage}
\renewcommand{\headrulewidth}{0.3pt}
\setlength{\parskip}{3pt}
\setlength{\parindent}{1.2em}
\setlength{\emergencystretch}{2em}
\setlist{itemsep=2pt,topsep=4pt}
\numberwithin{equation}{section}
\newtheorem{theorem}{Theorem}[section]
\newtheorem{lemma}[theorem]{Lemma}
\newtheorem{proposition}[theorem]{Proposition}
\newtheorem{corollary}[theorem]{Corollary}
\theoremstyle{definition}
\newtheorem{definition}[theorem]{Definition}
\newtheorem{example}[theorem]{Example}
\theoremstyle{remark}
\newtheorem{remark}[theorem]{Remark}
\newcommand{\mot}{\mathop{o}}
\newcommand{\W}{\mathcal W}
\newcommand{\J}{\mathcal J}
\newcommand{\F}{\mathcal F}
\newcommand{\G}{\mathcal G}
\newcommand{\eps}{\varepsilon}
\newcommand{\len}{\mathop{\rm len}}
\newcommand{\down}{\mathord{\downarrow}}
\newcommand{\emb}{\preccurlyeq}
\newcommand{\eqv}{\simeq}
\newcommand{\N}{\mathbb N}
\newcommand{\Hn}[1]{H_{#1}}
\newcommand{\Wordspace}[1]{s_{\omega^2}^{F}(#1)}
\newtcolorbox{resultbox}{enhanced,breakable,colback=pale,colframe=accent,
  boxrule=0.6pt,arc=1mm,left=3mm,right=3mm,top=2mm,bottom=2mm}
\lstset{basicstyle=\ttfamily\small,breaklines=true,columns=fullflexible,
  frame=single,rulecolor=\color{accent},backgroundcolor=\color{pale},
  showstringspaces=false,tabsize=4}
\IfFileExists{data/test_summary.tex}{\newcommand{\TestAssertions}{3,240,040}
\newcommand{\ProductChecks}{1,625,127}
\newcommand{\MarkerStages}{355}
\newcommand{\DeciderChecks}{774,149}
\newcommand{\BinaryPairs}{726,242}
\newcommand{\BinaryStages}{14}
}{
\newcommand{\TestAssertions}{1,710,857}
\newcommand{\ProductChecks}{1,580,460}
\newcommand{\MarkerStages}{355}
\newcommand{\DeciderChecks}{47,907}}

\title{\vspace{-10mm}\color{ink}\LARGE\bfseries
Exact maximal order types below $\omega^2$\\[3mm]
\large A proposed solution to Altman's finite-alphabet\\sharpness question}
\author{Research draft prepared with ChatGPT\\
\small For Vladimir Reshetnikov}
\date{September 19, 2026}

\begin{document}
\maketitle
\thispagestyle{empty}
\begin{abstract}
For a finite poset $P$, consider all words of ordinal length less than
$\omega^2$, ordered by increasing subsequence embeddings that may increase
labels in $P$. Write $n=|P|$ and let $j(P)$ count all downsets of $P$, including
the empty one. This article gives a complete proposed proof that, for $n\ge2$,
\[
  \mot\bigl(\Wordspace P\bigr)=\omega^{\omega^{\,n+j(P)-2}}.
\]
In particular, the two-letter antichain has maximal order type
$\omega^{\omega^4}$, answering the explicit sharpness question in
Example~3.15 of Altman's March 2026 revision. The two-letter chain has type
$\omega^{\omega^3}$. The main construction embeds every finite Cartesian
power of an earlier stage into the next stage by synchronizing periodic
$\omega$-blocks. Proper downsets admit an excluded-letter guard; the universal
block instead requires a limit-length padding argument. A more general
classification for selected families of periodic tails isolates the exact
exception in which naive block counting fails. An exact symbolic embedding
algorithm and a reproducible test suite accompany the proofs.
\end{abstract}

\begin{resultbox}
\textbf{Status and scope.} This is a new, unrefereed research argument, not a
report of an independently certified result. All proof steps beyond the
explicitly stated classical maximal-order-type theorems are supplied below.
The computer checks test symbolic embeddings, not infinite ordinal ranks.
The claim concerns finite alphabet posets and the strict cutoff $\omega^2$;
it does not settle arbitrary alphabets or higher cutoffs.
\end{resultbox}

\vfill
\noindent\small\textbf{Keywords.} Well-quasi-order; maximal order type;
transfinite word; Higman ordering; finite downset lattice; natural product.
\newpage
\begingroup
\small
\setlength{\parskip}{0pt}
\tableofcontents
\endgroup
\newpage

\section{The selected question and the result}
\label{sec:intro}

Altman studies finite-image words of bounded transfinite length in
\cite{Altman2026}. Example~3.15, on printed page~10 of version~2, gives the
upper bound $\omega^{\omega^4}$ for a two-element antichain and explicitly
leaves its sharpness unverified. The same example improves the chain's
upper bound to $\omega^{\omega^3}$. The arXiv record checked for this note
lists version~2, dated March~10, 2026, as the latest revision. Targeted
searches did not locate a subsequent resolution; this is a limited
literature check, not a guarantee of priority.

This is an attractive problem to attack because the alphabet is tiny, all
$\omega$-tails have finite descriptions up to embedding, and the candidate
ordinal is explicit. The difficult part is not finding more tail types. It
is proving that the types already available contribute their full ordinal
complexity after concatenation.

Let $P$ be a finite poset. A \emph{downset} is a subset $D\subseteq P$ such
that $x\le_P y\in D$ implies $x\in D$. Put
\[
 \J(P)=\{D\subseteq P:D\text{ is a downset}\},\qquad j(P)=|\J(P)|.
\]
The empty set is counted in $j(P)$ throughout.

\begin{theorem}[Finite-alphabet formula]\label{thm:main}
Let $P$ be a finite poset with $n\ge2$ elements. With increasing,
label-respecting subsequence embedding and mutual embeddability quotiented
out,
\begin{equation}\label{eq:main}
 \boxed{\mot\bigl(\Wordspace P\bigr)
       =\omega^{\omega^{\,n+j(P)-2}}.}
\end{equation}
For the empty alphabet the maximal order type is $1$, and for a singleton
alphabet it is $\omega^2$.
\end{theorem}

Here $s_{\omega^2}^{F}(P)$ denotes words of length \emph{strictly less than}
$\omega^2$. Since $P$ is finite, the finite-image superscript is automatic.
The theorem is proved in Section~\ref{sec:classification} after two product
embedding lemmas.

\begin{corollary}[The selected binary cases]\label{cor:binary}
For a two-element antichain $A_2$ and a two-element chain $C_2$,
\[
 \mot\bigl(\Wordspace {A_2}\bigr)=\omega^{\omega^4},
 \qquad
 \mot\bigl(\Wordspace {C_2}\bigr)=\omega^{\omega^3}.
\]
\end{corollary}
\begin{proof}
There are four downsets in $A_2$ and three in $C_2$. Substitute $n=2$ in
\eqref{eq:main}.
\end{proof}

The proof establishes more than these two values. We may choose an arbitrary
family of nonempty downsets and allow only the corresponding periodic
$\omega$-blocks. Unless the universal block is the sole allowed infinite
block, the number of allowed block types determines the maximal order type
exactly. The exceptional case is essential, not a technical inconvenience.

\subsection{What the proof adds}
The upper bound follows from a finite generating alphabet. The lower bound
requires an additional construction. A new periodic block serves as a
separator only when an embedding cannot hide it inside earlier blocks.
Introducing downsets in inclusion order solves that problem. A finite guard
then prevents a component from leaking into the next separator. At the
universal separator no excluded letter exists, so a different device is
needed: attach a cancellable tail of limit length. An infinite tail cannot
fit into a bounded, hence finite, initial segment of an $\omega$-block.

Together these observations give actual order embeddings
\[
  X^{r+1}\hookrightarrow Y\qquad(r<\omega)
\]
from each stage into the next. Classical natural-product arithmetic turns
these embeddings into the sharp lower bound.

\section{Definitions and classical inputs}
\label{sec:prelim}

\subsection{Ordinal words and embeddability}
A \emph{word over $P$} is a function $u:\alpha\to P$, where $\alpha$ is an
ordinal; its length is $\len(u)=\alpha$. If $u:\alpha\to P$ and
$v:\beta\to P$, write $u\emb v$ when there is a strictly increasing map
$f:\alpha\to\beta$ satisfying
\[
 u(\xi)\le_P v(f(\xi))\qquad(\xi<\alpha).
\]
No continuity or cofinality condition is imposed on $f$. Write $u\eqv v$
when $u\emb v$ and $v\emb u$. Concatenation is ordinal concatenation. The
empty word $\eps$ has length zero. For lengths below $\omega^2$ there are
unique $k,m<\omega$ such that
\[
 \len(u)=\omega k+m.
\]

Embeddability is in general only a quasi-order. Distinct words, such as
$a^\omega$ and $aa^\omega$, may be equivalent. All maximal order types in
this article refer to the quotient by $\eqv$.

\begin{lemma}[Concatenation is monotone]\label{lem:concat}
If $u_i\emb v_i$ for $0\le i<r$, where $r$ is finite, then
$u_0\cdots u_{r-1}\emb v_0\cdots v_{r-1}$.
\end{lemma}
\begin{proof}
Translate the embedding on each component into the corresponding interval
of the target concatenation. The intervals occur in the same order, so the
union of these translated maps is strictly increasing and respects labels.
\end{proof}

A finite concatenation of words of length below $\omega^2$ still has length
below $\omega^2$. In particular, all maps used later stay within the stated
cutoff, even though the number of $\omega$-blocks in their arguments is
unbounded across different arguments.

\subsection{Maximal order types}
A quasi-order is a \emph{well-quasi-order} if every infinite sequence
$x_0,x_1,\ldots$ contains $i<j$ with $x_i\le x_j$. The maximal order type
$\mot(X)$ is the largest ordinal type of a linear extension of the quotient
poset. The existence of this maximum, its characterization by the tree of
bad finite sequences, and the product theorem below are classical
\cite{deJonghParikh1977}. We use the following precise inputs.

\begin{theorem}[Classical maximal-order-type calculus]\label{thm:classical}
For well-quasi-orders $X,Y$:
\begin{enumerate}[label=\textnormal{(\roman*)}]
\item An order embedding $X\hookrightarrow Y$ implies $\mot(X)\le\mot(Y)$.
\item An order-preserving surjection from $Y$ onto $X$, including surjectivity
up to quasi-order equivalence, implies $\mot(X)\le\mot(Y)$.
\item For the componentwise Cartesian product,
$\mot(X\times Y)=\mot(X)\otimes\mot(Y)$, where $\otimes$ is natural product.
\item If $P$ is a nonempty finite poset of cardinality $n$, its Higman word
ordering satisfies
\begin{equation}\label{eq:finiteword}
 \mot(P^*)=H_n,\qquad H_n:=\omega^{\omega^{n-1}}.
\end{equation}
\end{enumerate}
\end{theorem}

The finite-word formula is part of the de Jongh--Parikh--Schmidt theory
\cite{deJonghParikh1977,Schmidt2020}. These are imported theorems; the present
article does not claim new proofs of them. All later lower bounds are
reduced to these inputs by explicit embeddings.

For clarity, assertion~(ii) has an elementary bad-sequence explanation.
Choose a preimage for each element of the target quotient. If a target
sequence is bad, its sequence of chosen preimages is bad, because any
comparison in the source would project to a comparison in the target.
This embeds the target bad-sequence tree into the source tree. The same
argument shows the target is a well-quasi-order whenever the source is.

\subsection{The ordinal arithmetic used here}
For ordinals in Cantor normal form, natural sum $\oplus$ adds coefficients
at equal exponents, and natural product $\otimes$ multiplies the resulting
formal polynomials, using natural sum on exponents. In particular,
\[
 (\omega^\alpha)\otimes(\omega^\beta)=\omega^{\alpha\oplus\beta}.
\]
Consequently, for finite $N\ge1$ and $q\ge1$,
\begin{equation}\label{eq:power}
 H_N^{\otimes q}=\omega^{\omega^{N-1}q},
\end{equation}
where the multiplication by $q$ on the right is ordinary ordinal
multiplication by a positive integer. Continuity of ordinal exponentiation
gives
\begin{equation}\label{eq:sup}
 \sup_{q\ge1}H_N^{\otimes q}
 =\omega^{\sup_{q\ge1}\omega^{N-1}q}
 =\omega^{\omega^N}=H_{N+1}.
\end{equation}
Equation~\eqref{eq:sup} is the entire ordinal calculation behind the
induction. Ordinary powers and natural powers must not be interchanged in
general; \eqref{eq:power} states exactly the special case being used.

\section{Periodic tails and the finite generating alphabet}
\label{sec:normalform}

\subsection{One periodic block for each nonempty downset}
For a nonempty downset $D=\{d_0,\ldots,d_{t-1}\}$, fix a periodic word
\begin{equation}\label{eq:block}
 B_D=(d_0\cdots d_{t-1})^\omega.
\end{equation}
The chosen enumeration is immaterial up to $\eqv$: every element of $D$
appears infinitely often in every such representative.

\begin{lemma}[Comparison of periodic blocks]\label{lem:blocks}
For nonempty downsets $D,E$ and a letter $a\in P$,
\[
 a\emb B_D\ \Longleftrightarrow\ a\in D,
 \qquad
 B_D\emb B_E\ \Longleftrightarrow\ D\subseteq E.
\]
Every nonempty tail of $B_D$ is equivalent to $B_D$.
\end{lemma}
\begin{proof}
If $a\le_P d$ for a letter $d\in D$, downward closure gives $a\in D$.
Conversely, $a\in D$ supplies a matching occurrence in $B_D$. If
$B_D\emb B_E$, applying the same observation to an occurrence of each
$d\in D$ gives $D\subseteq E$. If $D\subseteq E$, recursively match the
letters of $B_D$ to successive identical occurrences in $B_E$. There is
always a later occurrence, so the recursion gives an embedding.
The identical recursive construction works in any nonempty tail of a
periodic block; the reverse embedding into the full block is inclusion.
\end{proof}

For a nonempty subset $R\subseteq P$, repeating $R$ or repeating its downset
$\down R$ gives equivalent $\omega$-words. In the less immediate direction,
each $x\in\down R$ can be matched to an element $r\in R$ above it, and each
such $r$ recurs infinitely often.

\subsection{Reduction of arbitrary words}
The following periodic-tail reduction is standard; compare
\cite[Propositions~3.3--3.4]{Altman2026}. A proof is included to specify
exactly which equivalence is needed.

\begin{lemma}[Finite prefix and recurrent tail]\label{lem:recurrent}
Every word $u:\omega\to P$ is equivalent to $vB_D$ for some finite word
$v\in P^*$ and some nonempty downset $D$.
\end{lemma}
\begin{proof}
Let $R$ be the set of labels occurring infinitely often in $u$. Since $P$
is finite and $u$ is infinite, $R$ is nonempty. Every label outside $R$
occurs only finitely often. Choose an index after all such occurrences,
and write $u=vt$, where $v$ is finite and $t$ uses only labels from $R$.
Every element of $R$ occurs infinitely often in $t$.

Put $D=\down R$. To embed $t$ into $B_D$, successively choose later identical
occurrences of its labels. To embed $B_D$ into $t$, for each requested label
$d\in D$ choose an $r\in R$ with $d\le_P r$, and then choose an occurrence
of $r$ later than all previous choices. Infinitely many occurrences of $r$
make this possible. Thus $t\eqv B_D$, and monotonicity of concatenation
gives $u\eqv vB_D$.
\end{proof}

\begin{proposition}[Finite representation]\label{prop:representation}
Every word of length less than $\omega^2$ is equivalent to a finite
concatenation of one-letter words and blocks $B_D$, with
$\varnothing\ne D\in\J(P)$.
\end{proposition}
\begin{proof}
Split a word of length $\omega k+m$ into its $k$ consecutive intervals of
length $\omega$ and its final finite interval. Apply
Lemma~\ref{lem:recurrent} to each infinite interval and concatenate the
resulting representatives.
\end{proof}

This is a representation theorem, not a uniqueness theorem. Different
finite expressions can represent equivalent words. For example,
$aB_{\down a}\eqv B_{\down a}$. Such absorption is precisely why a
finite generating alphabet gives only an upper bound without further work.

\begin{definition}[Selected families]\label{def:selected}
For $\F\subseteq\J(P)\setminus\{\varnothing\}$, let $\W(P,\F)$ be the
quasi-order of all finite concatenations of letters in $P$ and periodic
blocks $B_D$ with $D\in\F$, using actual ordinal-word embeddability.
Finite expressions are identified up to mutual embeddability when maximal
order type is taken.
\end{definition}

In particular, $\W(P,\varnothing)=P^*$, and
$\W(P,\J(P)\setminus\{\varnothing\})$ has the same quotient order as
$\Wordspace P$ by Proposition~\ref{prop:representation}.

\begin{proposition}[Atom-count upper bound]\label{prop:upper}
For $n=|P|\ge1$ and $m=|\F|$,
\begin{equation}\label{eq:upper}
 \mot(\W(P,\F))\le H_{n+m}.
\end{equation}
Moreover, $\W(P,\F)$ is a well-quasi-order.
\end{proposition}
\begin{proof}
Take a discrete alphabet of $n+m$ symbols: one for each finite letter and
one for each selected periodic block. Evaluate a finite symbol string by
ordinal concatenation. If one string is a subsequence of another, the
corresponding evaluated word embeds into the other by
Lemma~\ref{lem:concat}. Evaluation is surjective onto the represented
quasi-order. Apply the finite-word formula and monotone-surjection
principle in Theorem~\ref{thm:classical}.
\end{proof}

Using a discrete alphabet deliberately ignores extra comparisons among the
atoms. This does not weaken the desired upper bound: every poset with the
same positive finite cardinality has the same finite-word maximal order
type. No order-reflection claim is made for the evaluation map.

\section{Cofinal localization and synchronization}
\label{sec:sync}

The following lemma concerns \emph{displayed target tokens}. A finite word
may contain many one-letter tokens adjacent to each other, and a token
expression need not be a canonical decomposition of its ordinal length.
Nevertheless, each displayed $B_E$ occupies a convex interval of order type
$\omega$, and there are only finitely many displayed tokens.

\begin{lemma}[Cofinal localization]\label{lem:cofinal}
Suppose a periodic block $B_D$ embeds into a finite token expression $v$.
Then a tail of its image lies in one displayed target block $B_E$, is
cofinal in that block, and satisfies $D\subseteq E$.
If a source position occurs after the whole source block $B_D$, its image
occurs after the whole target block $B_E$.
\end{lemma}
\begin{proof}
Number the target tokens in their order of occurrence. As the source
positions $0,1,2,\ldots$ increase, the indices of the tokens containing
their images form a nondecreasing sequence with only finitely many
possible values. It is eventually constant. The final token cannot be a
one-letter token, because the embedding is injective. It is therefore
some $B_E$. The image contains infinitely many positions of that token,
so it is unbounded in its order type $\omega$.

Every $d\in D$ still occurs in the source tail. Each such occurrence maps
to a label in $E$ above $d$. Downward closure gives $d\in E$, proving
$D\subseteq E$. Finally, the image of any subsequent source position must
lie above every image point of $B_D$. Since those points are cofinal in the
target $B_E$, no position of that target block remains available.
\end{proof}

\begin{remark}[A block may straddle earlier blocks]\label{rem:straddle}
The lemma does not say that the entire source block lands in a single
target block. An initial finite part can be placed earlier. Only the tail
is localized. All subsequent proofs use this weaker, correct statement.
\end{remark}

\begin{lemma}[Equal-count synchronization]\label{lem:sync}
Let $D$ be a nonempty downset, and suppose an old family $\G$ satisfies
\begin{equation}\label{eq:newness}
  D\nsubseteq E\qquad\text{for every }E\in\G.
\end{equation}
Consider two token expressions with exactly $r$ displayed copies of $B_D$,
all other infinite tokens belonging to $\G$. Under an embedding of the
first expression into the second, the tail of the $i$th displayed source
$B_D$ is cofinal in the $i$th displayed target $B_D$, for $1\le i\le r$.
\end{lemma}
\begin{proof}
By Lemma~\ref{lem:cofinal}, each displayed source $B_D$ requires a target
infinite block whose downset contains $D$. Condition~\eqref{eq:newness}
excludes all old target blocks. Its cofinal destination is therefore one
of the $r$ displayed new target blocks. Successive source new blocks have
strictly increasing cofinal destination indices: after the first has used
a target block cofinally, any later source position must be beyond it.
A strictly increasing injection from $\{1,\ldots,r\}$ into itself is the
identity. This proves the claim.
\end{proof}

To arrange \eqref{eq:newness}, introduce a finite collection of downsets in
an inclusion-respecting order. Ordering them by cardinality, with arbitrary
tie breaking, suffices: a new set cannot be contained in an earlier,
distinct set of no larger size. The full downset $P$, when present, is last.

\section{Proper markers: finite guards give product embeddings}
\label{sec:proper}

\begin{lemma}[Right cancellation of a final letter]\label{lem:cancel}
For arbitrary ordinal words $u,v$ and any fixed letter $c$,
\[
 uc\emb vc\quad\Longrightarrow\quad u\emb v.
\]
\end{lemma}
\begin{proof}
The image of the final source letter lies at or before the final target
letter. Every position belonging to $u$ maps strictly before that image,
hence inside the target prefix $v$. Restrict the embedding.
\end{proof}

\begin{theorem}[Proper-marker product embedding]\label{thm:proper}
Let $D$ be a nonempty proper downset of $P$, and let $\G$ satisfy
\eqref{eq:newness}. Choose $c\in P\setminus D$ and put
$X=\W(P,\G)$, $Y=\W(P,\G\cup\{D\})$. For every $r\ge1$, the map
\begin{equation}\label{eq:propermap}
 \Phi_r(u_0,\ldots,u_r)
   =u_0cB_Du_1cB_D\cdots u_{r-1}cB_Du_r
\end{equation}
is an order embedding $X^{r+1}\hookrightarrow Y$.
\end{theorem}
\begin{proof}
Order preservation follows by concatenating componentwise embeddings and
identity maps on the separators. It remains to prove reflection.

Suppose $\Phi_r(u_0,\ldots,u_r)\emb\Phi_r(v_0,\ldots,v_r)$.
By Lemma~\ref{lem:sync}, the $i$th source separator $B_D$ has image tail
cofinal in the $i$th target separator $B_D$. Thus, for $i>0$, every image
point of $u_i$ and of its subsequent guard lies after the $i$th target
separator. For $i=0$ there is no preceding separator and no lower-bound
restriction to impose.

Fix $i<r$. The source guard $c$ immediately following $u_i$ precedes the
next source separator. Some point of that next separator maps into the
corresponding next target separator. Therefore the source guard's image
is before the end of that target separator. It cannot be in that separator:
if it were matched to a label $d\in D$, then $c\le_P d$ would imply
$c\in D$, contrary to the choice of $c$. Its image consequently lies in
the target interval $v_ic$ before the separator. Every point of $u_i$
precedes the guard, so its image lies in $v_i$, not in the final guard of
$v_ic$. We obtain $u_i\emb v_i$.

After the last source separator, cofinality forces the entire source
suffix $u_r$ to map into the target suffix $v_r$. Hence $u_r\emb v_r$ as
well. All components compare in the product order. Reflection and
preservation imply an order embedding after taking equivalence classes.
\end{proof}

The guard is allowed to map into the interior of $v_i$; it need not match
the explicitly displayed target guard. The proof only needs the fact that
it lies no later than that guard. This is important when old components
already contain many occurrences of $c$ or old infinite blocks.

\begin{corollary}[One proper marker raises the level]\label{cor:properlevel}
Under the hypotheses of Theorem~\ref{thm:proper}, if
$\mot(X)=H_N$ and $n+|\G|=N$, then $\mot(Y)=H_{N+1}$.
\end{corollary}
\begin{proof}
For every $r\ge1$, product embedding and the classical product theorem give
$\mot(Y)\ge H_N^{\otimes(r+1)}$. Taking the supremum and using
\eqref{eq:sup} yields $\mot(Y)\ge H_{N+1}$. Proposition~\ref{prop:upper}
gives the reverse inequality.
\end{proof}

\section{The universal marker: limit padding replaces a guard}
\label{sec:universal}

For $D=P$ there is no excluded letter. Any fixed finite word $z$ is absorbed
by the universal periodic block:
\begin{equation}\label{eq:absorb}
 zB_P\eqv B_P.
\end{equation}
Indeed, the countable sequence $zB_P$ can be matched greedily into $B_P$;
every requested label reappears infinitely often. Thus attaching a finite
guard before $B_P$ cannot preserve arbitrary component information.
The remedy is to end the component with a cancellable \emph{infinite} tail.

\subsection{A cancellable map into limit lengths}

\begin{lemma}[Limit padding]\label{lem:padding}
Let $E$ be a nonempty proper downset of $P$, and choose $c\notin E$.
The map
\begin{equation}\label{eq:padding}
 G(u)=ucB_E
\end{equation}
is order preserving and order reflecting on arbitrary words of length
below $\omega^2$. Every $G(u)$ has nonzero limit length, of the form
$\omega k$ with $k\ge1$. If $B_E$ is allowed in a family, $G$ maps that
family's represented word order into itself.
\end{lemma}
\begin{proof}
Monotonicity is immediate. Suppose $ucB_E\emb vcB_E$. The source occurrence
of $c$ cannot map into the final target $B_E$, by downward closure and
$c\notin E$. Thus it maps into $vc$, and the restriction to $u$ maps into
$v$, by the same final-letter argument as Lemma~\ref{lem:cancel}. Hence
$u\emb v$.

If $\len(u)=\omega k+m$, then
\[
 \len(ucB_E)=(\omega k+m)+1+\omega=\omega(k+1),
\]
which is a nonzero limit ordinal. The family assertion follows from the
explicit concatenation.
\end{proof}

\begin{lemma}[No infinite tail fits below a finite bound]\label{lem:noleak}
Let $u$ have nonzero limit length below $\omega^2$. Suppose an embedding
maps all of $u$ to positions preceding a fixed position $\eta$ inside a
target interval of order type $\omega$. Then no position of $u$ can map
inside that interval.
\end{lemma}
\begin{proof}
If a source position $\xi$ mapped into the interval before $\eta$, there
would be only finitely many target positions between its image and
$\eta$. But a nonzero limit ordinal has infinitely many positions after
each of its positions: otherwise the ordinal would have a last position.
The infinitely many source positions after $\xi$ would have to inject into
that finite target interval. This is impossible.
\end{proof}

The position $\eta$ need not be the first position of the target
$\omega$-block. Any point of it gives a finite bound. Nor is it necessary
for the next source marker to begin inside that block: one image point of
its cofinal tail suffices.

\subsection{The product embedding}

\begin{theorem}[Universal-marker product embedding]\label{thm:full}
Let $\G$ consist of nonempty proper downsets of $P$, with
$\G\ne\varnothing$. Choose $E\in\G$ and the map $G$ from
Lemma~\ref{lem:padding}. Put
$X=\W(P,\G)$ and $Y=\W(P,\G\cup\{P\})$. For each $r\ge1$,
\begin{equation}\label{eq:fullmap}
 \Psi_r(u_0,\ldots,u_r)
  =G(u_0)B_PG(u_1)B_P\cdots G(u_{r-1})B_Pu_r
\end{equation}
is an order embedding $X^{r+1}\hookrightarrow Y$.
\end{theorem}
\begin{proof}
Preservation follows from the monotonicity of $G$ and concatenation.
Suppose the image of $(u_0,\ldots,u_r)$ embeds into the image of
$(v_0,\ldots,v_r)$.

No old downset contains $P$. Lemma~\ref{lem:sync} therefore applies to the
$r$ displayed universal separators. Their cofinal destinations are the
corresponding $r$ displayed target universal separators.

Fix $i<r$. When $i>0$, cofinality at the preceding separator forces all of
$G(u_i)$ to map after the corresponding preceding target separator.
Choose an image point $\eta$ of the following source $B_P$ inside the
corresponding following target $B_P$; such a point exists by cofinal
localization. Every image point of $G(u_i)$ is before $\eta$.
By Lemma~\ref{lem:padding}, $G(u_i)$ has nonzero limit length. Applying
Lemma~\ref{lem:noleak} shows that none of its positions can be in this
following target $B_P$. Thus the whole image of $G(u_i)$ lies in the
intervening target component $G(v_i)$. For $i=0$ the same argument applies
without a preceding boundary. We have proved
\[
 G(u_i)\emb G(v_i)\qquad(i<r).
\]
Reflection for $G$ yields $u_i\emb v_i$ for all $i<r$.

After the last source universal separator, cofinality forces the final
source component $u_r$ into the final target component $v_r$. Hence
$u_r\emb v_r$, completing product reflection.
\end{proof}

\begin{corollary}[One universal marker raises the level when padding exists]
\label{cor:fulllevel}
Under the hypotheses of Theorem~\ref{thm:full}, if
$\mot(X)=H_N$ and $n+|\G|=N$, then $\mot(Y)=H_{N+1}$.
\end{corollary}
\begin{proof}
Apply the finite Cartesian-product theorem to the embeddings $\Psi_r$,
take the supremum in \eqref{eq:sup}, and combine with the atom-count
upper bound, exactly as in Corollary~\ref{cor:properlevel}.
\end{proof}

Notice the role of the old proper block. It need not encode the whole
component by itself. Its job is to give a limit-length ending, while an
excluded finite letter makes that ending cancellable. These two properties
are logically distinct and both are used.

\section{The full classification and the main theorem}
\label{sec:classification}

\begin{theorem}[Selected-tail classification]\label{thm:family}
Let $P$ be a nonempty finite poset of size $n$, and let
$\F\subseteq\J(P)\setminus\{\varnothing\}$.
\begin{enumerate}[label=\textnormal{(\roman*)}]
\item If $\F\ne\{P\}$, then
\begin{equation}\label{eq:family}
 \mot(\W(P,\F))=H_{n+|\F|}
 =\omega^{\omega^{\,n+|\F|-1}}.
\end{equation}
This includes the empty family.
\item If $\F=\{P\}$, then
\begin{equation}\label{eq:exception}
 \mot(\W(P,\{P\}))=H_n\cdot\omega
    =\omega^{\omega^{n-1}+1}.
\end{equation}
\end{enumerate}
\end{theorem}

\begin{proof}[Proof of part~\textnormal{(i)}]
For the empty family, \eqref{eq:family} is the finite-word formula.
Otherwise $\F\ne\{P\}$ implies that the family contains at least one
proper downset. List its proper downsets as $D_1,\ldots,D_t$ in increasing
cardinality, breaking ties arbitrarily. If $P\in\F$, append it last.

Start with $X_0=P^*$, of type $H_n$. At step $i\le t$, no earlier downset
contains $D_i$. Indeed, an earlier downset has no greater cardinality, and
equality together with containment would make the downsets equal.
Theorem~\ref{thm:proper} applies. Inductively,
Corollary~\ref{cor:properlevel} gives
\[
 \mot\bigl(\W(P,\{D_1,\ldots,D_i\})\bigr)=H_{n+i}.
\]
If $P$ is absent from $\F$, this finishes the proof. If $P$ is present,
there is an old proper downset because $t\ge1$. Apply
Theorem~\ref{thm:full} and Corollary~\ref{cor:fulllevel} to the final step,
obtaining $H_{n+t+1}=H_{n+|\F|}$.
\end{proof}

\begin{proof}[Proof of part~\textnormal{(ii)}]
Every expression using only the universal infinite block can be written
\[
 u_0B_Pu_1B_P\cdots u_{k-1}B_Pv,
 \qquad u_i,v\in P^*.
\]
By \eqref{eq:absorb}, it is equivalent to $B_P^kv$. Here $B_P^k$ means
$k$ successive copies, not an infinite repetition. These normal forms
satisfy
\begin{equation}\label{eq:univorder}
 B_P^kv\emb B_P^\ell w
 \quad\Longleftrightarrow\quad
 \bigl(k<\ell\bigr)\ \text{or}\
 \bigl(k=\ell\text{ and }v\emb w\bigr).
\end{equation}
If $k<\ell$, send the $k$ source infinite blocks into the first $k$ target
blocks and the finite suffix into one extra universal target block.
If $k>\ell$, $k$ successive source infinite blocks require $k$ distinct
cofinal target infinite destinations by Lemma~\ref{lem:cofinal}, which is
impossible. If $k=\ell$, synchronization leaves the source suffix wholly
in the target suffix, so $v\emb w$ is necessary. It is plainly sufficient.

Thus the quotient is the ordered sum of $\omega$ copies of $P^*$:
\begin{equation}\label{eq:sum}
 \W(P,\{P\})/\eqv\ \cong\ \sum_{k<\omega}P^*.
\end{equation}
Every linear extension must put all of the $k$th copy before the
$(k+1)$st copy. Each copy has maximal order type $H_n$, and one may use
a maximal linear extension in every copy. Therefore the maximal order
type is the ordinary ordered sum $\sum_{k<\omega}H_n=H_n\cdot\omega$.
The elementary identity
$\omega^\alpha\cdot\omega=\omega^{\alpha+1}$ gives
\eqref{eq:exception}.
\end{proof}

\begin{proof}[Proof of Theorem~\ref{thm:main}]
Suppose $n\ge2$. Choose a minimal element $a$ of $P$. The singleton
$\{a\}$ is a nonempty proper downset. Consequently the full family
$\J(P)\setminus\{\varnothing\}$ is not the exceptional family $\{P\}$.
Its size is $j(P)-1$. Proposition~\ref{prop:representation} and
Theorem~\ref{thm:family}(i) give
\[
 \mot\bigl(\Wordspace P\bigr)
 =H_{n+j(P)-1}=\omega^{\omega^{n+j(P)-2}}.
\]
For $n=1$, there is one word at each ordinal length $\alpha<\omega^2$, and
embedding is exactly comparison of lengths. Its quotient is therefore the
ordinal $\omega^2$ itself. Equivalently, apply part~(ii) with $H_1=\omega$.
For $n=0$, the empty word is the sole word, giving maximal order type $1$.
\end{proof}

\subsection{Why the supremum argument is legitimate}
The proof never identifies the order type of an increasing union with the
supremum of the order types of its pieces. Such an identification is not a
general principle for maximal order types. Instead, for each fixed finite
$r$, there is a separately verified order embedding of a Cartesian power.
It gives a genuine lower bound on the \emph{same} target order type.
Taking the supremum of a collection of lower bounds is valid. The upper
bound is independently obtained by a monotone surjection.

\subsection{The singleton warning is structural}
For a singleton alphabet, two formal atoms are available: the letter and
its periodic $\omega$-block. Their finite words would have maximal type
$H_2=\omega^\omega$ if treated as a free Higman alphabet. But after ordinal
evaluation, the actual type is only $\omega^2$. The family
$\F=\{P\}$ exhibits the same defect over every nonempty finite alphabet.
Hence atom counting alone cannot be a lower-bound proof. The padding
construction explains exactly why the full family over $n\ge2$ avoids this
collapse.

\section{The binary proof written out explicitly}
\label{sec:binary}

This section specializes the general proof so the selected problem can be
checked without tracking an arbitrary downset family.

\subsection{Two incomparable letters}
Let $P=\{a,b\}$ with no comparisons between distinct letters, and abbreviate
\[
 A=a^\omega,\qquad B=b^\omega,\qquad U=(ab)^\omega.
\]
There are exactly three nonempty downsets, giving these three periodic
block types. Define stages by their allowed tokens:
\[
 \begin{array}{c|c|c}
  \text{stage}&\text{allowed tokens}&\text{claimed maximal type}\\\hline
  X_0&a,b&H_2=\omega^{\omega}\\
  X_1&a,b,A&H_3=\omega^{\omega^2}\\
  X_2&a,b,A,B&H_4=\omega^{\omega^3}\\
  X_3&a,b,A,B,U&H_5=\omega^{\omega^4}.
 \end{array}
\]
The upper bound in each row comes from the number of tokens.
The first equality is classical. To pass from $X_0$ to $X_1$, use, for
every finite $r$,
\[
 (u_0,\ldots,u_r)\longmapsto
 u_0bAu_1bA\cdots u_{r-1}bAu_r.
\]
Source $A$-blocks must synchronize with target $A$-blocks, and the guard
$b$ cannot enter $A$. This embeds $X_0^{r+1}$.

To pass from $X_1$ to $X_2$, use
\[
 (u_0,\ldots,u_r)\longmapsto
 u_0aBu_1aB\cdots u_{r-1}aBu_r.
\]
A source $B$-block cannot be cofinal in an old target $A$-block, and its
preceding guard $a$ cannot enter a target $B$-block. This embeds
$X_1^{r+1}$.

For the final step, set
\[
 G(u)=ubA.
\]
The guard $b$ cannot enter the final $A$, so $G$ reflects embedding; the
image has limit length. Now use
\[
 (u_0,\ldots,u_r)\longmapsto
 G(u_0)U G(u_1)U\cdots G(u_{r-1})Uu_r.
\]
A source $U$ contains both letters infinitely often, so neither $A$ nor
$B$ can be its cofinal destination. The displayed $U$-blocks synchronize.
Limit padding prevents the intervening components from entering the next
$U$ before a fixed image point of its source marker. This embeds
$X_2^{r+1}$.

At each of the three steps, \eqref{eq:sup} converts these product embeddings
into the next displayed ordinal. Finally, every binary word of length
below $\omega^2$ is equivalent to an expression in $a,b,A,B,U$.
This proves the sharpness of the binary antichain bound.

\subsection{Two ordered letters}
Let $0<1$. The nonempty downsets are $\{0\}$ and $\{0,1\}$; their periodic
blocks may be represented by
\[
 A=0^\omega,\qquad U=1^\omega.
\]
The latter is equivalent to $(01)^\omega$. Begin with finite words over
$\{0,1\}$, of maximal type $H_2$. Add $A$ using the guard $1$, giving
$H_3$. Then add $U$ using the limit padding $G(u)=u1A$, giving $H_4$.
Thus the chain's sharpened upper bound is also attained.

\subsection{Three easily confused comparisons}
For the antichain, the following facts illustrate the mechanism:
\[
 aU\eqv U,\qquad A\emb U,\qquad AB\not\emb U.
\]
The first is finite-prefix absorption. The second uses cofinal recurrence
of $a$ in $U$. The third follows because two successive $\omega$-blocks
cannot both have cofinal image in the single target block. Also
$U\not\emb AB$: its tail would have to lie in either the pure $a$ block
or the pure $b$ block, losing one of its recurrent labels. These facts
show why both order of occurrence and cofinal type matter.

\section{Consequences and finite-poset examples}
\label{sec:consequences}

\subsection{Chains, antichains, and extremal values}

\begin{corollary}\label{cor:extrema}
For $n\ge2$,
\[
 \mot\bigl(\Wordspace {C_n}\bigr)=\omega^{\omega^{2n-1}},
 \qquad
 \mot\bigl(\Wordspace {A_n}\bigr)=\omega^{\omega^{n+2^n-2}}.
\]
Among $n$-element posets, the first value is the minimum, attained exactly
by chains, and the second is the maximum, attained exactly by antichains.
\end{corollary}
\begin{proof}
A chain has exactly $n+1$ downsets, its initial segments. Every subset of
an antichain is a downset, so it has $2^n$ downsets. For any finite poset,
the initial segments of a linear extension give $n+1$ distinct downsets,
and there are at most $2^n$ subsets in total.

For the equality statement at the lower end, fix a linear extension. If
there are only $n+1$ downsets, every downset is one of its initial segments.
In particular all principal downsets are comparable by inclusion. Since
$\down x\subseteq\down y$ implies $x\le_P y$, every pair of elements is
comparable, and $P$ is a chain. At the upper end, if every subset is a
downset, a strict relation $x<_P y$ would make $\{y\}$ fail to be a downset.
Thus $P$ is an antichain. The function of $j(P)$ in \eqref{eq:main} is
strictly increasing, so the ordinal extremal assertions follow.
\end{proof}

\subsection{Equivalent counting invariants}
Every downset of a finite poset is generated by its maximal elements,
which form an antichain. Conversely, the maximal elements of the downset
generated by an antichain are that antichain. These maps give a bijection,
including the empty set. Thus $j(P)$ may equally be described as the number
of antichains of $P$.

Complementation gives a bijection between downsets and upsets. Passing to
the dual poset exchanges downsets with upsets, so $j(P^{\mathrm{op}})=j(P)$.
The formula therefore yields
\[
 \mot\bigl(\Wordspace {P^{\mathrm{op}}}\bigr)
 =\mot\bigl(\Wordspace P\bigr).
\]
This is equality of maximal order types, not a claim that the word
quasi-orders are isomorphic.

For a disjoint union $P\amalg Q$, choosing a downset is equivalent to
choosing one downset in each component, hence
$j(P\amalg Q)=j(P)j(Q)$. For an ordered sum $P+Q$, in which every element
of $P$ lies below every element of $Q$, one has
$j(P+Q)=j(P)+j(Q)-1$. The two possibilities are a downset contained in $P$,
or all of $P$ together with a nonempty downset of $Q$. These identities
make \eqref{eq:main} explicit for many recursively assembled alphabets.

\subsection{A table of examples}
In Table~\ref{tab:examples}, $\rho(P)=n+j(P)-2$ is the finite exponent
parameter in $\omega^{\omega^{\rho(P)}}$.

\begin{table}[htbp]
\centering
\begin{tabular}{@{}lrrrl@{}}
\toprule
Alphabet $P$ & $n$ & $j(P)$ & $\rho(P)$ & Maximal order type\\
\midrule
Two-element chain &2&3&3&$\omega^{\omega^3}$\\
Two-element antichain &2&4&4&$\omega^{\omega^4}$\\
Three-element chain &3&4&5&$\omega^{\omega^5}$\\
Three-element $V$ or its dual &3&5&6&$\omega^{\omega^6}$\\
Two-element chain plus an isolated point &3&6&7&$\omega^{\omega^7}$\\
Three-element antichain &3&8&9&$\omega^{\omega^9}$\\
Four-element chain &4&5&7&$\omega^{\omega^7}$\\
Four-element diamond &4&6&8&$\omega^{\omega^8}$\\
Four-element antichain &4&16&18&$\omega^{\omega^{18}}$\\
Five-element antichain &5&32&35&$\omega^{\omega^{35}}$\\
\bottomrule
\end{tabular}
\caption{Exact consequences of Theorem~\ref{thm:main}. The diamond has one
bottom element, two incomparable middle elements, and one top element.}
\label{tab:examples}
\end{table}

The alphabet cardinality by itself is not enough: the two binary values
already differ. Conversely, nonisomorphic alphabets can have the same
maximal order type. For instance the three-element $V$ and its dual have
the same $n$ and $j$, while a four-element chain and a three-element chain
plus an isolated point have the same value of $n+j-2$ despite different
cardinalities.

\subsection{Finite quasi-orders}
For a finite quasi-order rather than a poset, first pass to its quotient
by mutual comparison. Projection of labels preserves and reflects word
embedding: a comparison of quotient labels is exactly the comparison of
any representatives. Therefore the theorem applies with $n$ equal to the
number of equivalence classes and $j$ equal to the number of downsets of
the quotient poset. Counting raw labels instead of classes can give an
incorrect answer.

\section{An exact symbolic embedding algorithm}
\label{sec:algorithm}

\subsection{Representation}
The supplied module \texttt{code/omega2.py} represents a word by a finite
sequence of tokens. A token \texttt{('f', x)} is a finite letter, and
\texttt{('w', D)} is $B_D$, with $D$ stored as a nonempty bitmask downset.
The poset relation is supplied by predecessor masks. The code validates
reflexivity, antisymmetry, and transitivity.

This representation does not claim to compute the recurrent set of an
arbitrary black-box infinite input. Rather, it decides embeddings between
already supplied finite periodic-block descriptions. The mathematical
normal-form theorem says such descriptions represent every equivalence
class in the target order.

\subsection{Greedy decision procedure}
Maintain a pointer to a target token, initially the first. Read source
tokens from left to right.

\begin{description}[style=nextline,leftmargin=0pt,labelwidth=0pt,labelindent=0pt]
\item[Source token is a finite letter $a$.]
Scan target tokens until finding either a finite letter $b$ with $a\le_P b$,
or an infinite token $B_E$ with $a\in E$. In the finite case consume the
letter and advance the pointer. In the infinite case leave the pointer at
$B_E$: after finitely many matches, an equivalent periodic tail remains.
\item[Source token is an infinite block $B_D$.]
Scan target tokens until finding an infinite token $B_E$ with $D\subseteq E$.
Consume its tail cofinally and advance the pointer to the following token.
\end{description}
If the required token does not exist, reject. Accept after every source
token has been processed.

\begin{theorem}[Exactness of the algorithm]\label{thm:algorithm}
The greedy procedure accepts exactly when the represented source ordinal
word embeds into the represented target ordinal word.
\end{theorem}
\begin{proof}
Within a periodic target block, removing any finite initial segment leaves
an equivalent tail. Thus the symbolic state need not record a phase or a
finite offset in that block.

For a finite source letter, the first eligible target token provides an
early possible match. Choosing a later token cannot leave a more useful
suffix. Within an infinite token, choosing an actual sufficiently early
eligible occurrence leaves an equivalent periodic tail. Therefore a
successful continuation after any possible match can also be made after
the greedy match.

For an infinite source block $B_D$, Lemma~\ref{lem:cofinal} shows that every
embedding must use some target block $B_E$ cofinally with $D\subseteq E$.
All later source positions must then lie after that target block. The
greedy choice takes the first eligible such block. The entire source
$B_D$ embeds in its available tail by Lemma~\ref{lem:blocks}; allowing
part of the source block to occur earlier cannot improve the endpoint
beyond this first eligible block. The remaining target suffix after the
greedy choice contains the suffix left by any later choice.

Inducting on the number of remaining source tokens proves completeness:
whenever a completion exists, the greedy step preserves the existence of
a completion. Conversely, each accepted finite step can be realized by
an actual letter match in a finite token or a periodic tail, and each
accepted infinite step by a cofinal embedding into the selected target
block. Concatenating these realizations gives an ordinal embedding, so
acceptance is sound.
\end{proof}

The target pointer never retreats. With bit operations and relation lookup
counted as unit cost, the decision procedure uses
$O(|u|_{\mathrm{tok}}+|v|_{\mathrm{tok}})$ time and constant auxiliary space.
For unbounded alphabet size, bitmask operations have their ordinary bit
costs; validating masks and the poset incurs additional cost. The
reference checker uses memoized exploration of every eligible target
endpoint rather than the first endpoint.

\subsection{Why finite expansion is not an adequate substitute}
Replacing each $B_D$ by $M$ repetitions for a large integer $M$ changes the
problem. The absorption equivalence $aB_{\down a}\eqv B_{\down a}$ already
fails for the direction from length $M+1$ to length $M$ after a uniform
finite expansion. Conversely, finite approximations do not directly
certify the cofinal-boundary restrictions needed for reflection of product
maps. The supplied code therefore treats $\omega$-blocks symbolically and
uses the proven cofinal semantics exactly.

\section{Computational checks and reproducibility}
\label{sec:tests}

The included verification run passed \TestAssertions\ assertions. Its
random seed is \texttt{20260919}; all random checks supplement deterministic
checks. The report is stored in \texttt{data/verification.json}. The
recorded run used Python~3.13.5, and the code targets Python~3.10 or newer
without third-party packages.

\begin{table}[htbp]
\centering
\begin{tabular}{@{}lr@{}}
\toprule
Check & Number of assertions\\
\midrule
Exhaustive greedy/reference agreement & 38,107\\
Seeded longer-word greedy/reference agreement & 9,800\\
Exhaustive proper/universal product maps & 1,544,960\\
Seeded longer-word product maps & 35,500\\
Padding reflection & 55,150\\
Padding has nonzero limit length & 4,024\\
Universal-only comparison formula & 20,640\\
Universal-only normalization & 2,660\\
Boundary, absorption, and regression cases & 16\\
\midrule
Total & \TestAssertions\\
\bottomrule
\end{tabular}
\caption{Executed symbolic checks. A ``word length'' in the test generator
means token length, not ordinal length.}
\label{tab:checks}
\end{table}

The exhaustive comparison of the two deciders uses all naturally labelled
posets of sizes zero through three and every source/target token word of
length at most two. The longer-word comparisons cover all naturally
labelled posets of sizes two through four and use token words of length
at most twelve.

For the product embeddings, every naturally labelled poset of sizes two
through four is processed, introducing all nonempty downsets in cardinality
order. There are \MarkerStages\ marker stages. At each stage the exhaustive
test uses a fixed four-word component pool and all pairs of tuples of
arities two and three. The seeded tests add tuples of arities two through
five with component token lengths at most eight. In total these give
\ProductChecks\ product-map comparisons. Separate padding tests allow all
one-token words, including already universal words, so padding reflection
is checked beyond the narrower family needed in the main induction.

The two embedding implementations share the mathematical periodic-tail
semantics. Agreement between them is a check of the greedy optimization,
not an independent derivation of that semantics. Likewise, no collection
of the finite assertions in Table~\ref{tab:checks} proves
$\omega^{\omega^4}$ to be the maximal order type. That conclusion rests on
the unbounded family of product embeddings and the ordinal calculation.

\subsection{Enumeration conventions}
The auxiliary enumerator considers relations on $\{0,\ldots,n-1\}$ whose
strict comparisons respect the usual numerical order. Every finite-poset
isomorphism class has such a labelling, but the enumeration is neither an
enumeration of all labellings nor an isomorphism-free enumeration. The
numbers of generated relations for $n=0,1,2,3,4,5$ are respectively
\[
 1,\ 1,\ 2,\ 7,\ 40,\ 357.
\]
These counts were computed by the included program. Histograms by the
number of downsets appear in \texttt{data/ideal\_counts.csv}. Table examples
are also exported to \texttt{data/examples.csv}. No claim about uniform
sampling of isomorphism classes is intended.

To compute $j(P)$ for one small input poset, the code simply checks all
$2^n$ subsets for downward closure. This exponential enumeration is
sufficient for the examples and the proof tests. The article does not
claim a polynomial-time algorithm for downset counting.

\subsection{Reproduce the package}
From the root of the unpacked directory:
\begin{lstlisting}
python3 code/verify.py
latexmk -pdf -interaction=nonstopmode -halt-on-error article.tex
\end{lstlisting}
The Python command regenerates the verification reports and CSV tables.
The PDF can alternatively be built by running \texttt{pdflatex article.tex}
repeatedly until cross-references stabilize. No external bibliography
processor is required; a separate BibTeX file is supplied for reuse.

\section{Proof audit, scope, and further directions}
\label{sec:audit}

\subsection{Where a hidden gap would have to occur}
The proof is organized so its nonclassical burden is concentrated in a few
claims. The atom-count upper bound is a monotone-surjection argument.
Cofinal localization uses finiteness of the target token decomposition.
Synchronization uses the newness condition $D\nsubseteq E$ for all old
blocks, not merely distinctness of block labels. The proper-marker
reflection uses downward closure to exclude a guard from the next marker.
The full-marker reflection uses \emph{both} cancellation and the limit
length of the padded component.

These dependencies explain several tempting but incorrect shortcuts.
A source infinite block need not land wholly inside one target token.
The ordinary evaluation map from atom strings need not reflect order.
The universal marker admits no excluded finite letter. A finite test
cannot certify an infinite supremum. None of these shortcuts is used.

An independent referee can focus first on Theorems~\ref{thm:proper} and
\ref{thm:full}. Once their product-reflection assertions are accepted, the
classification follows by a short finite induction and the classical
finite-product formula. The binary specialization in
Section~\ref{sec:binary} provides a second presentation of the same core
argument, not an independent proof with different foundations.

\subsection{Established within the argument versus not claimed}
Within this manuscript, the main theorem has a complete proof relative to
the classical results stated in Theorem~\ref{thm:classical}. There is no
remaining conjectural lemma inside that derivation. Nevertheless the
argument has not undergone independent expert review or proof-assistant
verification. The appropriate external status is therefore a
\emph{proposed solution}, with explicit hypotheses and reproducible
supporting checks, rather than a certified resolution.

The article does not compute all invariants of these quasi-orders: their
height, width, and exact isomorphism classification are different tasks.
It does not show that two alphabets with equal $n+j(P)$ produce isomorphic
word orders. It does not give a canonical representation for arbitrary
infinite input streams. It also does not assert that selected block
families at higher ordinal cutoffs have the same counting law.

\subsection{Why the proof stops at this cutoff}
The decisive compactness feature is that a target word below $\omega^2$
has only finitely many $\omega$-blocks. The image of a source $\omega$-block
must eventually stabilize in one of them. At length $\omega^2$, a source
$\omega$-block may instead visit infinitely many successive target
$\omega$-blocks, taking only finitely many points from each. Thus the
cofinal localization lemma in its present form is no longer valid when
that target length is allowed.

A higher-cutoff extension would need a hierarchy of separator ranks and a
new localization theorem that accounts for escape through infinitely many
lower-rank blocks. Counting indecomposable blocks without that theorem
would reproduce the same logical mistake that naive atom counting makes
at the universal-only stage.

For infinite alphabets, even when every individual word has finite image,
the set of possible recurrent downsets need not be finite, and a single
finite inclusion ordering does not exhaust the construction. The present
finite induction therefore does not directly transfer. These observations
identify precise missing mechanisms for extensions; they are not reasons
to reject further investigation.

\subsection{A concrete formal-verification target}
A useful formalization would first define ordinal concatenation and its
embedding relation for finite periodic-token codes, then prove the exact
symbolic decision theorem. The next targets are the cofinal boundary lemma,
proper-marker reflection, and universal-marker reflection. With the
classical maximal-order-type product and finite-word theorems available,
the final induction has very little additional complexity. The package
contains executable tests but no Lean development, and it makes no claim
of machine-checked proof.

\appendix
\clearpage
\section{A compact proof-dependency map}
\label{app:dependencies}

\begin{longtable}{@{}>{\raggedright\arraybackslash}p{0.29\textwidth}>{\raggedright\arraybackslash}p{0.31\textwidth}>{\raggedright\arraybackslash}p{0.32\textwidth}@{}}
\toprule
Statement & Inputs & Role\\
\midrule
\endhead
Finite representation, Proposition~\ref{prop:representation}
 & Finiteness of $P$; recurrent labels; monotone concatenation
 & Identifies the full word quotient with a finite-token quotient.\\[3pt]
Atom-count upper bound, Proposition~\ref{prop:upper}
 & Finite-word formula; monotone surjection
 & Gives $\mot(\W(P,\F))\le H_{n+|\F|}$.\\[3pt]
Cofinal localization, Lemma~\ref{lem:cofinal}
 & Finitely many target tokens; strict increase; periodic recurrence
 & Finds the target block used cofinally by each source marker.\\[3pt]
Synchronization, Lemma~\ref{lem:sync}
 & Cofinal localization; no old block contains the new downset
 & Forces the $i$th new marker to correspond to the $i$th new marker.\\[3pt]
Proper product embedding, Theorem~\ref{thm:proper}
 & Synchronization; a letter outside the proper downset
 & Embeds every finite Cartesian power of the old stage.\\[3pt]
Limit padding, Lemma~\ref{lem:padding}
 & A proper old block; an excluded letter; final-letter cancellation
 & Produces an order-reflecting map into nonzero limit lengths.\\[3pt]
Universal product embedding, Theorem~\ref{thm:full}
 & Synchronization; limit padding; the finite-bound contradiction
 & Adds the universal block without losing the product lower bound.\\[3pt]
Family classification, Theorem~\ref{thm:family}
 & The two product maps; classical natural product; atom upper bound
 & Computes every selected family, including its sole exception.\\[3pt]
Main theorem, Theorem~\ref{thm:main}
 & Full family has a proper downset when $n\ge2$; finite representation
 & Gives $\omega^{\omega^{n+j(P)-2}}$ and the two boundary cases.\\
\bottomrule
\end{longtable}

\section{Minimal worked use of the code}
\label{app:code}

The following example can be run with \texttt{code/} on the Python import
path. Downset masks \texttt{1}, \texttt{2}, and \texttt{3} represent
$\{a\}$, $\{b\}$, and $\{a,b\}$.

\begin{lstlisting}[language=Python]
from omega2 import FinitePoset, letter, omega, embeds, theorem_value

P = FinitePoset.antichain(2)
a, b = letter(0), letter(1)
A, B, U = omega(1), omega(2), omega(3)

assert embeds(P, (a, U), (U,))       # finite-prefix absorption
assert embeds(P, (A,), (U,))        # pure tail into universal tail
assert not embeds(P, (A, B), (U,))  # two cofinal blocks cannot share one
assert not embeds(P, (U,), (A, B))  # no eligible cofinal destination

print(theorem_value(P))
# omega^(omega^4)

C = FinitePoset.chain(2)
print(theorem_value(C))
# omega^(omega^3)

print(theorem_value(P, family=(P.full,)))
# omega^(omega^1 + 1) = omega^(omega + 1)
\end{lstlisting}

The last output illustrates the exceptional universal-only family. Its
maximal order type is much smaller than $H_3=\omega^{\omega^2}$, even
though the two finite letters and the universal block form three distinct
formal tokens.

\section{Literature and reproducibility record}
The chosen source was checked both in arXiv HTML and in the PDF rendering;
Example~3.15 is on printed page~10, which is zero-based PDF page~9. Its
version identifier is \texttt{2409.03199v2}. The source of the open question,
the two standard maximal-order-type references, and the distinction between
proven classical inputs and the present proposed argument are recorded in
\texttt{notes/literature\_audit.md}.

The archive contains this source and its compiled PDF, a reusable
\texttt{references.bib}, the symbolic module, the reproducible verification
script, machine-readable test results, example and downset-count tables,
a proof-audit note, and a README. It contains no third-party article PDFs,
font files, checksum files, or LaTeX build intermediates.

\begin{thebibliography}{9}
\bibitem{Altman2026}
Harry Altman.
\newblock \emph{Bounding finite-image sequences of length $\omega^k$}.
\newblock arXiv:2409.03199v2, revised March 10, 2026; manuscript dated
March 9, 2026. In particular, Example 3.15, p.~10.
\newblock \url{https://arxiv.org/abs/2409.03199v2}.

\bibitem{deJonghParikh1977}
D.~H.~J. de Jongh and Rohit Parikh.
\newblock Well-partial orderings and hierarchies.
\newblock \emph{Indagationes Mathematicae (Proceedings)}, 80(3):195--207,
1977; also cited as \emph{Indagationes Mathematicae} 39.
\newblock \url{https://doi.org/10.1016/1385-7258(77)90067-1}.

\bibitem{Schmidt2020}
Diana Schmidt.
\newblock Well-Partial Orderings and their Maximal Order Types.
\newblock In P.~Schuster, M.~Seisenberger, and A.~Weiermann, editors,
\emph{Well-Quasi Orders in Computation, Logic, Language and Reasoning},
Trends in Logic 53, pp.~351--391. Springer, 2020.
\newblock Reprint of the 1979 work.
\newblock \url{https://doi.org/10.1007/978-3-030-30229-0_13}.
\end{thebibliography}

\end{document}
